\textbf{Data.} Images are generated with numpy: a background of three smoothed-noise channels plus a sinusoidal grating, 3--5 objects of random class, size, position and orientation (circle; rotated square; equilateral triangle), each filled with a random colour plus its own smooth texture, pasted in random order so that objects overlap and occlude, and Gaussian pixel noise. A fixed pool of \rhval{const/pool_size} training images (generator seed \rhval{const/pool_seed}) and a fixed test set of \rhval{const/test_size} images (seed \rhval{const/test_seed}) are used by every run. A run with size $N$ draws $N$ images from the pool with a permutation seeded by the run seed, so seed $s$ at size $N$ always sees the same subset and noise realisation across loss functions; noise is applied once, before training, with a seed that depends on the run seed.

\textbf{Training.} AdamW with peak learning rate \rhval{const/learning_rate:3} (OneCycle schedule), weight decay \rhval{const/weight_decay:4}, batch size \rhval{const/batch_size}, \rhval{const/train_steps} optimiser steps for \emph{every} $N$ (so small $N$ is trained for many more epochs than large $N$; this fixed-compute choice is a design decision of the study and is probed in the training-length sweep), random horizontal flip, final checkpoint, no early stopping and no validation set (a clean validation set would not exist in the noisy-label setting). All runs use two CPU threads on a shared machine; a run takes roughly one to two minutes.

\textbf{Experiments.} \emph{main}: $N=500$, four systems (CE, GCE, SCE, band-ignore CE) $\times$ three tasks (clean, boundary noise $p=0.3$, object-flip noise $p=0.3$), \rhval{const/seeds_main} seeds (0--3). Band-ignore CE additionally ran under flip noise $p=0.1$ and $p=0.2$ (three seeds), so that the matching CE sweep has a reference. \emph{size}: CE only, $N\in\{100,250,1000,2000\}$; clean and flip tasks at all sizes and boundary noise at $N=100,2000$, four seeds at $N=100$ and $N=2000$ and three seeds at $N=250,1000$ (the $N=500$ column is the CE row of main). \emph{sweep\_noise}: CE, $N=500$, flip $p\in\{0.1,0.2\}$, three seeds. \emph{sweep\_gce\_q}: GCE with $q\in\{0.3,0.9\}$ under flip $p=0.3$, $N=500$, three seeds ($q=0.7$ is the four-seed main row). \emph{sweep\_steps}: CE at $N=100$ under flip $p=0.3$ with 200 and 1200 steps, three seeds (the 600-step row is in size). Seeds fix the data subset, the noise, the initialisation and the batch order; every metric is reproducible from them. The seed count is below the five suggested in the task description because of the runtime budget; the study has no boundary-noise levels other than $p=0.3$ and no noise-level sweep for boundary noise.

\textbf{Pilots (disclosed).} Two earlier pilots (\rhval{const/pilot1_steps} steps at learning rate \rhval{const/pilot_lr:3}, then \rhval{const/pilot_steps} steps at the same rate) left several seeds with a collapsed triangle or square class even on clean labels, so they could not show noise effects. Their rows are logged and marked superseded; the final budget (\rhval{const/train_steps} steps, learning rate \rhval{const/learning_rate:2}) was chosen from them and from two single-seed clean-label runs that were scored on the test split. This is a small test-split leak in the choice of one global optimiser budget; it is the same for every system and noise condition, and no per-system or per-noise tuning was done. GCE and SCE use the defaults of the original papers.

\textbf{Hardware and runtime.} Shared Apple-silicon CPU, two threads per run, two runs at a time; one run takes roughly one to two minutes and the whole study (including pilots) about two hours of wall-clock time.

\textbf{Statistics.} Verdicts use Welch $t$-tests ($\alpha=0.05$, no multiplicity correction) between per-seed values, as registered in \texttt{proposal.md}; for H4 the tests are those printed by \texttt{rh compare} against CE. With 3--5 seeds these tests have little power and we report standard deviations over seeds. All statistics come from logged runs.
